\startchapter[title={Pruebas funcionales}, reference=Anexo_Pruebas]

El Cuadro~\in[tab:pruebas_funcionales] lista los casos de prueba funcionales
descritos en la Sección~\in[sec:pruebas_funcionales]. Cuando un caso agrupa
varias pruebas, cada una se distingue con una letra.

\startplacetable[reference=tab:pruebas_funcionales, location={here,split}, title={Pruebas funcionales y qué verifica cada una.}]
	\bTABLE[split=yes, align=middle, frame=on, width=fit]
	\setupTABLE[row][1][background=color, backgroundcolor=darkgray, style=bold]
	\setupTABLE[column][2][width=13cm, align=flushleft]
	\bTABLEhead
	\bTR \bTH Caso \eTH \bTH Qué se prueba \eTH \eTR
	\eTABLEhead
	\bTABLEbody
	\bTR \bTD 01 \eTD \bTD Caso básico: 2 pedidos en 2 lotes con 2 {\em pickers}. Cada {\em picker} ejecuta un lote y el conjunto de arcos usados es el esperado. \eTD \eTR
	\bTR \bTD 02 \eTD \bTD Dos pasillos, un pedido en cada uno: cada lote sigue la ruta de su pasillo y cada {\em picker} ejecuta un lote. \eTD \eTR
	\bTR \bTD 03 \eTD \bTD La instancia del caso 02, con controles adicionales: cada {\em picker} ejecuta exactamente un lote, la carga está equilibrada y no hay superposición en el tiempo. \eTD \eTR
	\bTR \bTD 04 \eTD \bTD Dos pasillos con pocas ubicaciones requeridas por pedido: las rutas cubren las ubicaciones requeridas. \eTD \eTR
	\bTR \bTD 05 \eTD \bTD Dos pasillos en los que los pedidos comparten una ubicación: rutas esperadas, carga equilibrada entre {\em pickers} y sin superposición en el tiempo. \eTD \eTR
	\bTR \bTD 06a \eTD \bTD Dos rutas posibles que llegan a tiempo: cada lote tiene una ruta y termina en su fecha de entrega. \eTD \eTR
	\bTR \bTD 06b \eTD \bTD Dos rutas posibles, la larga llega tarde: se elige la corta. \eTD \eTR
	\bTR \bTD 07 \eTD \bTD Un lote que solo puede completarse repitiendo una ubicación: se espera infactibilidad. \eTD \eTR
	\bTR \bTD 08 \eTD \bTD La instancia del caso 07 con una ubicación adicional que evita la repetición: tiene solución y la ruta es la esperada. \eTD \eTR
	\bTR \bTD 09 \eTD \bTD Dos pedidos de un mismo lote en dos ubicaciones: ruta esperada. \eTD \eTR
	\bTR \bTD 10 \eTD \bTD Dos pedidos en lotes distintos: ruta esperada para cada lote, carga equilibrada y sin superposición. \eTD \eTR
	\bTR \bTD 11 \eTD \bTD Cuatro pedidos en dos lotes que comparten ubicaciones: rutas esperadas. \eTD \eTR
	\bTR \bTD 12 \eTD \bTD Sin pruebas activas: verificaba la capacidad del carro dentro del modelo, que ya no la contempla. \eTD \eTR
	\bTR \bTD 13 \eTD \bTD 12 pedidos en 6 lotes con 3 {\em pickers}, sobre una grilla de 4×4: carga equilibrada y sin superposición en el tiempo. \eTD \eTR
	\bTR \bTD 14 \eTD \bTD La misma instancia y los mismos controles que el caso 13. \eTD \eTR
	\bTR \bTD 15 \eTD \bTD 3 pedidos con varias ubicaciones cada uno, en un lote con un {\em picker}, sobre una grilla de 4×4: tiene solución. \eTD \eTR
	\bTR \bTD 16 \eTD \bTD 6 pedidos en 3 lotes con 2 {\em pickers}, sobre una grilla de 4×4 con pasillos transversales; cada pedido visita más de una ubicación: tiene solución. \eTD \eTR
	\bTR \bTD 17 \eTD \bTD Un lote en un grafo ramificado: la ruta es la esperada. \eTD \eTR
	\bTR \bTD 18 \eTD \bTD Un lote de 5 pedidos cuya ruta requiere un desvío: tiene solución. \eTD \eTR
	\bTR \bTD 19a \eTD \bTD Ruta extendida (un lote de 6 pedidos, 9 ubicaciones) con dos pedidos en la misma ubicación: tiene solución. \eTD \eTR
	\bTR \bTD 19b \eTD \bTD La misma ruta con fecha de entrega 0 y 100: con 0 hay demora y con 100 no. \eTD \eTR
	\bTR \bTD 19c \eTD \bTD Ruta extendida: la ruta es un camino válido que cubre todas las ubicaciones requeridas. \eTD \eTR
	\bTR \bTD 20 \eTD \bTD Pedidos con varias ubicaciones sobre una grilla: cada lote visita todas sus ubicaciones. \eTD \eTR
	\bTR \bTD 21 \eTD \bTD Una ubicación requerida queda inalcanzable (puente cortado): se espera infactibilidad. \eTD \eTR
	\bTR \bTD 22a \eTD \bTD Un lote que termina antes de su fecha de entrega: no hay demora. \eTD \eTR
	\bTR \bTD 22b \eTD \bTD Un lote que termina después de su fecha de entrega: no hay adelanto. \eTD \eTR
	\bTR \bTD 23a \eTD \bTD Dos lotes de distinta longitud de ruta y dos {\em pickers} de distinta velocidad: el lote de ruta larga va al {\em picker} rápido. \eTD \eTR
	\bTR \bTD 23b \eTD \bTD Los tiempos de ruta calculados son los que usa la asignación, sin superposición entre lotes de un mismo {\em picker}. \eTD \eTR
	\bTR \bTD 24a \eTD \bTD La cadena de frío excluye al {\em picker} lento: el lote se asigna al rápido. \eTD \eTR
	\bTR \bTD 24b \eTD \bTD Ningún {\em picker} cumple la cadena de frío: el chequeo previo no encuentra pares elegibles. \eTD \eTR
	\bTR \bTD 25a \eTD \bTD Grilla de 5×3 con 2 {\em pickers} y 4 lotes, uno de los cuales supera la capacidad del carro: se espera infactibilidad. \eTD \eTR
	\bTR \bTD 25b \eTD \bTD La misma grilla, con dos pedidos de un lote en la misma ubicación: tiene solución. \eTD \eTR
	\bTR \bTD 25c \eTD \bTD La misma grilla, con pedidos de varias ubicaciones: cada lote visita todas sus ubicaciones. \eTD \eTR
	\bTR \bTD 25d \eTD \bTD La misma grilla, con todos los lotes por encima de la capacidad del carro: se espera infactibilidad. \eTD \eTR
	\bTR \bTD 25e \eTD \bTD La misma grilla, con fechas de entrega extremas (0 y 100): tiene solución en ambos casos. \eTD \eTR
	\bTR \bTD 25f \eTD \bTD La misma grilla, con la carga de cada lote exactamente igual a la capacidad del carro: tiene solución. \eTD \eTR
	\bTR \bTD 25g \eTD \bTD La misma grilla, con lotes distintos que comparten ubicaciones: tiene solución. \eTD \eTR
	\bTR \bTD 26a \eTD \bTD Grafo con un puente cortado, 2 {\em pickers} y 4 lotes: una ubicación requerida queda inalcanzable; se espera infactibilidad. \eTD \eTR
	\bTR \bTD 26b \eTD \bTD El mismo grafo con el puente, con fechas de entrega muy holgadas y dos tamaños de pedido: tiene solución. \eTD \eTR
	\bTR \bTD 26c \eTD \bTD El mismo grafo, con lotes que comparten ubicaciones: tiene solución. \eTD \eTR
	\bTR \bTD 26d \eTD \bTD El mismo grafo, con una fecha de entrega ajustada y un tiempo de recolección por unidad bajo y alto: tiene solución en ambos casos. \eTD \eTR
	\bTR \bTD 27 \eTD \bTD La instancia del caso 01 con otros pesos en la función objetivo. \eTD \eTR
	\bTR \bTD 28 \eTD \bTD Fechas de entrega ajustadas y holgadas: ambas tienen solución, porque la demora absorbe la fecha ajustada. \eTD \eTR
	\bTR \bTD 29 \eTD \bTD Un lote vacío entre otros: tiene solución. \eTD \eTR
	\bTR \bTD 30a \eTD \bTD Capacidades de pedido nulas y fraccionarias dentro de la capacidad del carro: tiene solución. \eTD \eTR
	\bTR \bTD 30b \eTD \bTD Capacidades de pedido nulas y fraccionarias que superan la capacidad del carro: se espera infactibilidad. \eTD \eTR
	\bTR \bTD 31a \eTD \bTD Un pedido asignado a más de un lote: se rechaza como infactible. \eTD \eTR
	\bTR \bTD 31b \eTD \bTD Cada pedido asignado a un único lote: tiene solución. \eTD \eTR
	\bTR \bTD 32 \eTD \bTD Un pedido que por sí solo supera la capacidad del carro: se espera infactibilidad. \eTD \eTR
	\bTR \bTD 33a \eTD \bTD Tiempos de servicio distintos entre ubicaciones: tiene solución. \eTD \eTR
	\bTR \bTD 33b \eTD \bTD Tiempo de servicio igual a 1 en todas las ubicaciones: tiene solución. \eTD \eTR
	\bTR \bTD 33c \eTD \bTD Tiempo de servicio nulo: tiene solución. \eTD \eTR
	\bTR \bTD 34a \eTD \bTD El tiempo de ejecución de cada lote coincide con su parte fija más el tiempo de ruta dividido por la velocidad del {\em picker}. \eTD \eTR
	\bTR \bTD 34b \eTD \bTD Ningún {\em picker} cumple la cadena de frío: el modelo informa infactibilidad sin fallar. \eTD \eTR
	\bTR \bTD 34c \eTD \bTD Un {\em picker} cuyo tiempo de ejecución es exactamente el límite de frío sigue siendo elegible. \eTD \eTR
	\bTR \bTD 35a \eTD \bTD El límite de frío se aplica a cada lote y no al turno: un {\em picker} puede encadenar dos lotes aunque el segundo termine después del límite. \eTD \eTR
	\bTR \bTD 35b \eTD \bTD Un {\em picker} que supera el límite de frío por menos que la tolerancia numérica sigue siendo elegible. \eTD \eTR
	\bTR \bTD 35c \eTD \bTD La parte fija del lote cuenta para el límite de frío: ningún {\em picker} es elegible y se espera infactibilidad. \eTD \eTR
	\bTR \bTD 36a \eTD \bTD Con un {\em picker} el doble de rápido, conviene repartir los dos lotes. \eTD \eTR
	\bTR \bTD 36b \eTD \bTD Con un {\em picker} cuatro veces más rápido, conviene que ejecute los dos lotes. \eTD \eTR
	\bTR \bTD 36c \eTD \bTD Las velocidades nulas o negativas se rechazan antes de resolver. \eTD \eTR
	\bTR \bTD 36d \eTD \bTD La demora se calcula con el tiempo de ruta escalado por la velocidad del {\em picker}. \eTD \eTR
	\bTR \bTD 37a \eTD \bTD Con un peso de finalización menor que la cantidad de pedidos del lote, el lote espera hasta la fecha de entrega. \eTD \eTR
	\bTR \bTD 37b \eTD \bTD Con un peso de finalización mayor, el lote termina lo antes posible. \eTD \eTR
	\bTR \bTD 38 \eTD \bTD Pedidos en dos pasillos sin salida, que obligan a pasar más de una vez por el mismo nodo: se espera infactibilidad. \eTD \eTR
	\eTABLEbody
	\eTABLE
\stopplacetable

\stopchapter
